\begin{abstract}
Generative models now synthesize posters, advertisements, visual text, and structured presentations. Yet visually plausible output need not be editable, and a structured output need not be reliably edited. Existing reviews often organize research by task or architecture, making it difficult to compare systems delivering pixels, layout boxes, vector paths, layered assets, or executable design code. This survey develops a representation--editability--evaluation (R--E--V) crosswalk for generative graphic design. We distinguish delivered artifacts from supported downstream operations, and format-level affordances from editing tasks whose performance is actually measured. We synthesize representative work in layout generation, content-aware composition, typography, vector graphics, layered reconstruction, and agentic design. We identify open evaluation challenges in instruction adherence, preservation of non-target content, structured artifact validity, cross-format comparability, and human workflow efficiency. The search protocol and paper-level verification remain in progress, so this manuscript should not be interpreted as an exhaustive systematic review.
\end{abstract}
